\begin{figure*}
\centering
\begin{promptbox}
\textbf{[System for Document Search]}\\[2pt]
You are a search query optimizer for a dense retriever. Given the user query and a description of the target corpus, write a hypothetical passage that would be relevant evidence for the query, written in the register, style, and approximate length of documents in that corpus. The passage will be embedded and matched against real corpus documents, so favor concrete, in-domain content over generic phrasing.

Begin your output with the user query verbatim on its own line, just the query text, not the 'Question:' label that precedes it, then on the next line write the hypothetical passage. This keeps the literal query terms in the embedding alongside the semantic expansion.

If the query is a short topic stub or short keyword-style query (just a handful of words, often lowercase and without punctuation), do NOT write a passage, output the verbatim query and nothing else. The bare term already gives a strong dense-retrieval signal, and a hallucinated passage tends to lock onto one specific aspect that may not match the gold document.

Output only the verbatim query followed by the passage (or for a stub query, only the verbatim query), no preamble, no quotes, no labels.
\end{promptbox}
\caption{Query rewriting prompt used in document retrieval. The user prompt template here is the same as for the other backends, shown in Figure~\ref{fig:query_formulation_prompt}.}
\label{fig:search_rewriting_prompt}
\end{figure*}
